%% kapa: exam=1 round=21 year=2010 period=1 form=A subject=civil_law subject_ko=민법 seq=5 number=5
%% source: Q-Net 감정평가사 기출문제 (cst00309) artlSeq=5024181 file=extracted/1008/5024181/2010년 제21회 감정평가사 1차시험 기출문제/1교시.hwp
%% conversion: kordoc 4.24.51 -> tools/gen1.py (auto); review=proofread-sonnet; flags=none
\begin{question}{5}
물권적 청구권에 관한 설명으로 옳지 않은 것은? (다툼이 있으면 판례에 의함)
\begin{choices}
  \item 물권적 청구권은 물권과 분리하여 양도할 수 없다.
  \item 부동산이 양도담보권자의 명의로 등기가 되어 있더라도 양도담보권설정자는 그 부동산의 불법점유자에 대하여 불법점유의 배제를 청구할 수 있다.
  \item 명의신탁이 유효한 경우, 신탁자는 명의신탁약정을 해지하고 수탁자 명의의 등기말소를 청구할 수 있다.
  \item 미등기 무허가건물의 양수인도 그 건물의 불법점유자에 대하여 직접 자신의 소유권에 기하여 건물의 반환을 청구할 수 있다.
  \item 대항력 없는 임차인은 임대인의 제3자에 대한 물권적 청구권을 대위행사할 수 있다.
\end{choices}
\end{question}
\answer{4}
